\documentclass[12pt]{article}
\usepackage{amsmath}
\usepackage{graphicx}
\usepackage{booktabs}
\usepackage{hyperref}

\title{Notes on ``Towards Real-Time Calibration-Free LIBS Supported by Machine Learning”}
\author{}
\date{}

\begin{document}

\maketitle

\section{Introduction}
Laser-Induced Breakdown Spectroscopy (LIBS) is a spectroscopic technique used to determine elemental composition of materials by analyzing emission from laser-induced plasma. A major challenge in LIBS is performing \textbf{quantitative analysis without calibration samples}. This approach is known as \textbf{Calibration-Free LIBS (CF-LIBS)}.

Traditional CF-LIBS requires strong expertise because it involves plasma diagnostics such as:

\begin{itemize}
\item Electron density ($n_e$) estimation from Stark broadening
\item Temperature ($T$) estimation using Boltzmann or Saha-Boltzmann plots
\item Spectral line integration and deconvolution
\end{itemize}

The paper proposes using \textbf{deep learning} to directly infer plasma parameters and elemental composition from spectra, enabling near \textbf{real-time CF-LIBS analysis}.

\section{Dataset Generation}
Instead of using experimental spectra, the authors generated a large synthetic dataset using a radiative transfer simulation code called \textbf{MERLIN}. This simulator models plasma emission under Local Thermodynamic Equilibrium (LTE).

The spectral radiance is modeled using the radiative transfer equation:

\begin{equation}
L_{\lambda}(x_i, n_e, T) = L_{\lambda,T}^{0}\left(1-e^{-\alpha_{\lambda}(x_i,n_e,T),\ell_c}\right)
\end{equation}

where:

\begin{center}
\begin{tabular}{ll}
\toprule
Symbol & Meaning \\
\midrule
$L_{\lambda}$ & Spectral radiance \\
$x_i$ & Mole fraction of element $i$ \\
$n_e$ & Electron density \\
$T$ & Plasma temperature \\
$\alpha_{\lambda}$ & Spectral absorption coefficient \\
$\ell_c$ & Plasma core dimension \\
\bottomrule
\end{tabular}
\end{center}

A total of \textbf{$10^6$ synthetic spectra} were generated.

\section{Elements Considered}
The dataset includes mixtures of nine elements:

\begin{itemize}
\item Ag
\item Ar
\item B
\item C
\item Cu
\item H
\item Mo
\item Ni
\item W
\end{itemize}

Their mole fractions satisfy the compositional constraint:

\begin{equation}
\sum_{i=1}^{9} x_i = 1
\end{equation}

The compositions are sampled using a \textbf{Dirichlet distribution} so that all possible mixtures are represented.

Minimum detectable concentration:

\begin{equation}
x_i = 10^{-6}
\end{equation}

which corresponds to parts-per-million level detection.

\section{Spectral Range}
The simulated spectra cover the range:

\begin{equation}
200 \text{ nm} \leq \lambda \leq 800 \text{ nm}
\end{equation}

However, each sample only contains a \textbf{60 nm window}. The center wavelength $\lambda_c$ is randomly selected from:

\begin{equation}
\lambda_c \sim U(230,770)
\end{equation}

This mimics real experimental spectrometers which only observe part of the spectrum.

\section{Plasma Parameters}
Two plasma parameters are varied during simulation.

\subsection{Electron Density}

\begin{equation}
5\times10^{22} \leq n_e \leq 5\times10^{23} ; m^{-3}
\end{equation}

This parameter controls Stark broadening of spectral lines.

\subsection{Temperature}

\begin{equation}
10,000K \leq T \leq 13,000K
\end{equation}

Temperature determines the population of excited states through the Boltzmann distribution.

\section{Data Normalization}
Spectral intensity ranges from $10^9$ to $10^{15}$, therefore normalization is required.

First, logarithmic scaling is applied:

\begin{equation}
L = \log(L_{\lambda})
\end{equation}

Then Z-score normalization:

\begin{equation}
L^* = \frac{L - \bar{L}}{\sigma_L}
\end{equation}

where $\bar{L}$ is the mean and $\sigma_L$ the standard deviation.

\section{Compositional Data Transformation}
To improve learning, mole fractions are transformed into log-ratio space:

\begin{equation}
x_i^* = \log(x_i) - \frac{1}{9} \sum_{j=1}^{9} \log(x_j)
\end{equation}

This transformation removes the sum-to-one constraint and improves neural network stability.

\section{Neural Network Architecture}
A \textbf{Convolutional Neural Network (CNN)} with residual connections is used.

Input:

\begin{itemize}
\item Spectral vector with 10240 values
\item Represents range 200–800 nm
\item Only 60 nm window contains real data
\end{itemize}

The architecture has two parts:

\subsection{Feature Extraction}
Three convolutional stages extract spectral features such as:

\begin{itemize}
\item spectral lines
\item line shapes
\item line overlaps
\end{itemize}

\subsection{Regression Layers}
Three fully connected layers predict physical quantities.

Outputs:

\begin{itemize}
\item 9 elemental mole fractions
\item electron density $n_e$
\item temperature $T$
\item center wavelength $\lambda_c$
\end{itemize}

The model contains approximately $6 \times 10^7$ parameters.

\section{Training Configuration}

Dataset size:

\begin{itemize}
\item Total spectra: $10^6$
\item Used for training: $10^5$
\end{itemize}

Split:

\begin{center}
\begin{tabular}{lc}
\toprule
Dataset & Size \\
\midrule
Training & 90,000 \\
Testing & 10,000 \\
\bottomrule
\end{tabular}
\end{center}

Training parameters:

\begin{itemize}
\item Batch size: 32
\item Learning rate: $5 \times 10^{-5}$
\item Epochs: 100
\end{itemize}

\section{Loss Function}
The loss function includes a confidence weighting:

\begin{equation}
L(p,t,c)=\sum_i c_i^2 (p_i - t_i)^2 + \sum_i (1-c_i)^2
\end{equation}

where:

\begin{itemize}
\item $p_i$ = predicted value
\item $t_i$ = target value
\item $c_i$ = confidence score
\end{itemize}

The network therefore learns both prediction and its reliability.

\section{Plasma Parameter Prediction Results}
The model predicts plasma parameters with high accuracy.

\begin{center}
\begin{tabular}{lcc}
\toprule
Parameter & Correlation & Mean Error \\
\midrule
Electron density & 0.996 & $<5\%$ \\
Temperature & 0.996 & $<1\%$ \\
\bottomrule
\end{tabular}
\end{center}

This shows neural networks can replace traditional diagnostics methods.

\section{Element Quantification}
The paper compares two cases.

\subsection{Nickel (Ni)}
Nickel has many spectral lines across the full spectral range.

Prediction performance:

\begin{itemize}
\item Pearson correlation: 0.965
\item Mean error: about 26%
\end{itemize}

\subsection{Hydrogen (H)}
Hydrogen has only a few spectral lines (Balmer series):

\begin{itemize}
\item $H_\delta$ 410 nm
\item $H_\gamma$ 434 nm
\item $H_\beta$ 486 nm
\item $H_\alpha$ 656 nm
\end{itemize}

Prediction accuracy improves dramatically if the spectral window contains $H_\alpha$.

Filtered results:

\begin{itemize}
\item Pearson correlation: 0.984
\item Mean error: 18%
\end{itemize}

\section{Real-Time Performance}
The entire pipeline includes:

\begin{enumerate}
\item reading camera image
\item spectral calibration
\item spatial integration
\item neural network inference
\end{enumerate}

Average processing time:

\begin{equation}
0.084 ; s
\end{equation}

This corresponds to approximately:

\begin{equation}
10 ; Hz
\end{equation}

which matches the repetition rate of typical LIBS lasers.

\section{Conclusion}
The study demonstrates that deep learning can enable real-time calibration-free LIBS analysis. Using synthetic spectra generated by radiative transfer simulations, a CNN model can simultaneously predict:

\begin{itemize}
\item plasma temperature
\item electron density
\item multi-element composition
\end{itemize}

This approach eliminates the need for manual spectral analysis and opens the possibility of fast, automated chemical diagnostics.

\end{document}
